% ===== PRÉAMBULE CANONIQUE — à coller VERBATIM en tête de CHAQUE .tex =====
\documentclass[11pt,a4paper]{article}
\usepackage[utf8]{inputenc}\usepackage[T1]{fontenc}\usepackage{lmodern}
\usepackage[french]{babel}
\usepackage{amsmath,amssymb,mathtools}\usepackage{icomma}
\usepackage[version=4]{mhchem}          % \ce{...} : équations & formules
\usepackage{siunitx}\sisetup{locale=FR} % \qty{}{}, \si{}, \ang{}
\usepackage{chemfig}                    % \chemfig{...} : structures développées
\usepackage{xcolor}\usepackage{colortbl}
\usepackage{tikz}\usepackage{pgfplots}\pgfplotsset{compat=1.18}
\usepackage[most]{tcolorbox}\usepackage{enumitem}\usepackage{array,booktabs}
\usepackage[a4paper,margin=2.2cm]{geometry}
\usetikzlibrary{arrows.meta,calc,angles,quotes,positioning,patterns,backgrounds,shapes.geometric}
\definecolor{primary}{RGB}{222,49,99}\definecolor{primarydark}{RGB}{150,22,55}
\definecolor{defcol}{RGB}{0,114,178}\definecolor{methcol}{RGB}{52,92,148}
\definecolor{excol}{RGB}{0,135,90}\definecolor{attncol}{RGB}{200,40,55}
\tcbset{boxbase/.style={enhanced,breakable,boxrule=0.9pt,arc=2.5pt,
  left=3mm,right=3mm,top=2.3mm,bottom=2.3mm,fonttitle=\bfseries,coltitle=white}}
% --- boîtes TUTORIEL (noms EXACTS, ne pas renommer) ---
\newtcolorbox{cle}[1][]{boxbase,colback=defcol!6,colframe=defcol,
  title={\textsc{À retenir}\ifx\relax#1\relax\relax\else\textmd{ : \itshape #1}\fi}}
\newtcolorbox{methode}[1][]{boxbase,colback=methcol!7,colframe=methcol,sharp corners,
  title={\textsc{Méthode}\ifx\relax#1\relax\relax\else\textmd{ --- \itshape #1}\fi}}
\newtcolorbox{exemplebox}[1][]{boxbase,colback=excol!5,colframe=excol,
  title={\textsc{Exemple}\ifx\relax#1\relax\relax\else\textmd{ : \itshape #1}\fi}}
\newtcolorbox{piege}[1][]{boxbase,colback=attncol!6,colframe=attncol,
  title={\textsc{Piège}\ifx\relax#1\relax\relax\else\textmd{ : \itshape #1}\fi}}
\newtcolorbox{remarque}[1][]{boxbase,colback=black!4,colframe=black!45,
  title={\textsc{Remarque}\ifx\relax#1\relax\relax\else\textmd{ : \itshape #1}\fi}}
% --- boîtes EXERCICE / CORRIGÉ (noms EXACTS, auto-comptées) ---
\newtcolorbox[auto counter]{exo}[1][]{boxbase,colback=primary!4,colframe=primary,
  title={\textsc{Exercice}~\thetcbcounter\ifx\relax#1\relax\relax\else\textmd{ --- \itshape #1}\fi}}
\newtcolorbox[auto counter]{corrige}[1][]{boxbase,colback=excol!4,colframe=excol,
  title={\textsc{Corrigé}~\thetcbcounter\ifx\relax#1\relax\relax\else\textmd{ --- \itshape #1}\fi}}
\newcommand{\R}{\mathbb{R}}\newcommand{\N}{\mathbb{N}}\newcommand{\Z}{\mathbb{Z}}\newcommand{\ds}{\displaystyle}
% (un \usepackage{fancyhdr}+en-tête est possible ; il est ignoré côté web)
% =========================================================================
\begin{document}
\sisetup{per-mode=symbol}
\begin{center}\small\itshape On rappelle : $\log 2 \approx 0{,}30$,\quad $\log 3 \approx 0{,}48$.\end{center}

\begin{exo}[mélange de deux acides forts]
On mélange \qty{50}{mL} de \ce{HCl} à \qty{0,10}{mol\per L} et \qty{50}{mL} de \ce{HCl} à \qty{0,30}{mol\per L}
(\ce{HCl} est un acide fort).
\begin{enumerate}[label=\alph*)]
  \item Calcule le nombre total de moles d'ions \ce{H3O+}.
  \item En déduis $[\ce{H3O+}]$ dans le mélange, puis le pH.
\end{enumerate}
\end{exo}

\begin{exo}[base forte dibasique]
L'hydroxyde de baryum \ce{Ba(OH)2} est une base forte. On en prépare une solution à
$C = \qty{5,0e-3}{mol\per L}$.
\begin{enumerate}[label=\alph*)]
  \item Écris sa dissociation et donne $[\ce{OH-}]$.
  \item Calcule le pOH, puis le pH.
\end{enumerate}
\end{exo}

\begin{exo}[dilution d'un acide fort]
Une solution d'acide chlorhydrique a $\mathrm{pH} = 2$.
\begin{enumerate}[label=\alph*)]
  \item Quelle est $[\ce{H3O+}]$ ?
  \item On dilue cette solution \textbf{100 fois}. Calcule le nouveau pH.
  \item Énonce la règle générale reliant le facteur de dilution et la variation de pH.
\end{enumerate}
\end{exo}

\begin{exo}[acide faible : calculer et vérifier]
Un acide faible \ce{HA} de $\mathrm{p}K_a = 4{,}0$ est en solution à $C_a = \qty{0,10}{mol\per L}$.
\begin{enumerate}[label=\alph*)]
  \item Calcule le pH.
  \item Calcule le degré d'ionisation $\alpha$ (Ostwald) et vérifie que l'approximation $\alpha \ll 1$ est
        acceptable.
\end{enumerate}
\end{exo}

\begin{exo}[reconnaître une base faible]
Une solution d'ammoniac \ce{NH3} à $\qty{1,0e-3}{mol\per L}$ a un pH \textbf{mesuré} de $10{,}15$. On donne
$\mathrm{p}K_a(\ce{NH4+/NH3}) = 9{,}3$.
\begin{enumerate}[label=\alph*)]
  \item Calcule le pH \textbf{théorique} attendu si \ce{NH3} se comporte comme une base faible.
  \item Compare à la mesure : \ce{NH3} est-elle une base forte ou une base faible ?
\end{enumerate}
\end{exo}

\end{document}
